\section{Results}
\label{sec:results}

We generated $N{=}100$ synthetic longitudinal patient profiles using SA and compared
them against the $K{=}23$ real patients across four progressively stringent validation
levels. As a baseline, we also generated $N{=}100$ patients from a regularized
multivariate normal (MVN) distribution fitted to the same data.

\subsection{Marginal Plausibility}
\label{sec:results-l1}

We first assessed whether SA-generated patients reproduced per-feature summary
statistics across all three visits. The median mean relative error (MRE) across all
72 features was approximately 2\%, with the majority of features below 5\%
(Table~\ref{tab:summary}), indicating that SA faithfully captured the central
tendencies of the real population. Representative coagulation factors---including
Factor~II (MRE $<$ 1.5\%), Factor~VIII (MRE $<$ 3\%), antithrombin (MRE $<$ 2\%),
and fibrinogen (MRE $<$ 1.2\%)---showed particularly strong agreement across all
visits. MVN also performed well at this level (median MRE $\approx$ 2--3\%), which
was expected since moment-matching is guaranteed by construction.

\begin{table}[ht]
\centering
\caption{Per-visit mean relative error (MRE) between real and SA-generated synthetic
patients for representative coagulation features, computed as
$|\bar{x}_{\text{synth}} - \bar{x}_{\text{real}}| / |\bar{x}_{\text{real}}|$.}
\label{tab:summary}
\begin{tabular}{lccc}
\toprule
Feature & Visit 1 & Visit 2 & Visit 3 \\
\midrule
Factor II (\%) & 0.009 & 0.012 & 0.015 \\
Factor VIII (\%) & 0.017 & 0.030 & 0.025 \\
AT (\%) & 0.013 & 0.018 & 0.014 \\
Fibrinogen (mg/dL) & $<$0.001 & 0.008 & 0.012 \\
TF Peak (nM) & 0.022 & 0.019 & 0.027 \\
TF ETP (nM$\cdot$min) & 0.015 & 0.021 & 0.018 \\
\bottomrule
\end{tabular}
\end{table}

We next examined whether known physiological relationships were preserved in the
synthetic cohort. We plotted pairwise scatter plots of coagulation factor levels
against thrombin generation parameters across all visits and observed that
SA-generated patients occupied the same joint regions as real patients
(Fig.~\ref{fig:bio-correlations}). The expected inverse relationship between
antithrombin and TF-initiated peak thrombin was preserved, as were the positive
associations between Factor~VIII and peak, prothrombin and ETP, and the inverse
relationship between antithrombin and ETP. These results indicated that SA captured
not only univariate distributions but also the biologically meaningful pairwise
dependencies that underlie coagulation physiology.

We also evaluated whether the characteristic pregnancy-driven longitudinal
trajectories were reproduced. We computed per-visit means and standard deviations
for six key coagulation features and found that synthetic patients tracked the
real trajectories closely, with overlapping confidence bands at all three visits
(Fig.~\ref{fig:pregnancy-progression}). Fibrinogen, Factor~VIII, and von Willebrand
factor showed the expected monotonic increases from baseline through the third
trimester in both cohorts, while Factor~II increased modestly and antithrombin
declined then partially recovered, consistent with the known hypercoagulable shift
of pregnancy. These patterns confirmed that the concatenated-visit SA pipeline
preserved the directionality and magnitude of longitudinal change.

\begin{figure}[ht]
\centering
\includegraphics[width=\textwidth]{validate_bio_correlations.pdf}
\caption{\textbf{Physiological correlations are preserved in synthetic patients.}
Scatter plots of coagulation factor levels versus thrombin generation parameters
across all visits (marker shapes indicate visit). Real patients (blue) and
SA-generated synthetic patients (tan) occupy the same joint regions, confirming
that biologically meaningful pairwise dependencies---including the AT--peak
inverse relationship (top left), the FVIII--peak positive relationship (top
right), and the FII--ETP and AT--ETP associations (bottom)---are preserved.}
\label{fig:bio-correlations}
\end{figure}

\begin{figure}[ht]
\centering
\includegraphics[width=\textwidth]{validate_pregnancy_progression.pdf}
\caption{\textbf{Pregnancy-driven longitudinal trajectories are faithfully
reproduced.} Mean $\pm$ standard deviation bands for six key coagulation
features across visits (BL = baseline, 1st = first trimester, 3rd = third
trimester) for real (blue) and SA-generated synthetic (orange) patients.
The characteristic pregnancy-driven increases in fibrinogen, Factor~VIII,
and vWF are captured, as are the stable-to-declining patterns in Factor~II
and AT.}
\label{fig:pregnancy-progression}
\end{figure}

\subsection{Cross-Visit Covariance Structure}
\label{sec:results-l2}

We compared the full $216 \times 216$ cross-visit correlation matrices for real,
SA-generated, and MVN-generated populations. The real data exhibited a characteristic
block structure with strong within-visit correlations along the diagonal and
structured off-diagonal blocks reflecting cross-visit dependencies---for example,
a patient's Visit~1 Factor~X level predicting their Visit~3 Factor~X level. We found
that SA preserved this block structure, including the off-diagonal cross-visit
correlations (Fig.~\ref{fig:cross-visit-corr}). The SA$-$Real residual matrix showed
small, unstructured deviations distributed throughout, whereas the MVN$-$Real residual
was notably smoother, reflecting the regularization-induced shrinkage of off-diagonal
correlations toward zero. This difference was most apparent in the off-diagonal blocks,
where MVN systematically underestimated cross-visit dependencies that SA retained.

We examined the eigenvalue spectrum of the sample covariance matrix to understand the
structural origin of this difference. The real data had rank~22 (eigenvalues dropping
to numerical zero at component~23), reflecting the $K{=}23$ patient constraint. SA's
spectrum dropped at component~18 (the PCA truncation point), faithfully reflecting
the low-rank structure without introducing variance in null dimensions. In contrast,
MVN's Ledoit--Wolf regularization inflated eigenvalues beyond component~22, producing
spurious variance in 194 dimensions where the data contained none (Supplementary
Fig.~S3). This spurious variance manifested as the increased dispersion visible in
PCA projections: SA-generated patients occupied the same region of PCA space as real
patients across all three visits, while MVN-generated patients showed substantially
greater scatter, particularly at Visits~2 and~3 (Fig.~\ref{fig:pca-by-visit}).

\begin{figure}[ht]
\centering
\includegraphics[width=\textwidth]{validate_cross_visit_corr.pdf}
\caption{\textbf{Cross-visit correlation structure.} Top row: full $216 \times 216$
Pearson correlation matrices for Real ($K{=}23$, left), SA ($N{=}100$, center),
and MVN ($N{=}100$, right) populations. Each matrix is organized as a
$3 \times 3$ grid of $72 \times 72$ blocks (delineated by black lines), where
the diagonal blocks capture within-visit feature correlations (V1--V1, V2--V2,
V3--V3) and the off-diagonal blocks capture cross-visit dependencies (e.g.,
V1--V3: does a patient's baseline Factor~VIII predict their third-trimester
Factor~VIII?). SA preserves both the within-visit and cross-visit block
structure. Bottom row: pairwise residual matrices (element-wise subtraction),
all plotted on the same $\pm 1$ color scale as the top row. The SA$-$Real
residual shows small, unstructured deviations distributed throughout. The
MVN$-$Real residual is notably smoother---particularly in the off-diagonal
blocks---reflecting the Ledoit--Wolf regularization that systematically shrinks
cross-visit correlations toward zero. The SA$-$MVN residual confirms that the
two methods differ most in precisely these off-diagonal blocks, where SA
retains longitudinal dependencies that MVN's rank-deficient covariance estimate
cannot represent. A version with amplified residual color scale ($\pm 0.5$) and
Frobenius norms is provided in Supplementary Fig.~S1.}
\label{fig:cross-visit-corr}
\end{figure}

\begin{figure}[ht]
\centering
\includegraphics[width=\textwidth]{sa_vs_mvn_pca_by_visit.pdf}
\caption{\textbf{PCA projections by visit: SA vs.\ MVN.} Each panel shows the
first two principal components (PC1: 27.7\% variance, PC2: 15.4\%) computed
from standardized per-visit real patient data ($K{=}23$ patients per visit,
72 features). Dark markers indicate real patients; lighter markers indicate
synthetic patients ($N{=}100$). Top row: SA-generated patients (circles)
cluster tightly around the real data cloud at all three visits, indicating
that SA preserves the geometry of the data manifold. Bottom row: MVN-generated
patients (diamonds) show substantially greater dispersion, particularly at
Visits~2 (first trimester) and~3 (third trimester), where MVN patients extend
well beyond the convex hull of the real data. This excess scatter reflects the
spurious variance introduced by Ledoit--Wolf regularization of the
rank-deficient covariance (Supplementary Fig.~S3). Colors encode visit: teal
(V1, not pregnant), green (V2, first trimester), orange (V3, third trimester).}
\label{fig:pca-by-visit}
\end{figure}

\subsection{Conditional Generation of Rare Subgroups}
\label{sec:results-l3}

We tested whether SA could amplify rare clinical subpopulations while preserving their
condition-specific signatures. Our dataset contained three clinical subgroups---Healthy
($n{=}14$), PCOS ($n{=}3$), and PE ($n{=}5$)---of which the PCOS and PE groups were
too small for any independent statistical analysis. We used SA's multiplicity-weighted
sampling to generate condition-specific cohorts of 100 patients each by upweighting
attention on the relevant stored patterns.

We compared condition-specific feature means between real and SA-generated patients
across eight coagulation features that exhibited between-condition variation
(Fig.~\ref{fig:conditioned-features}). SA-generated cohorts preserved the
condition-specific patterns: PCOS patients showed elevated Factor~VIII and vWF
relative to Healthy patients in both real and synthetic data, PE patients showed
elevated $\alpha_2$-AP and Factor~IX, and the between-condition rank ordering was
maintained across all eight features. To quantify equivalence, we performed a
bootstrap Mann--Whitney test: for each feature and condition, we subsampled the
synthetic cohort to match the real sample size ($n_{\text{real}}$), applied a
Mann--Whitney U test, and repeated 1,000 times. We reported the fraction of
replicates where $p > 0.05$ (i.e., real and synthetic were statistically
indistinguishable). Across all 24 feature--condition pairs, 19 (79\%) achieved
$p > 0.05$ in $\geq$90\% of replicates, with a median non-significance fraction
of 96.2\%. The five features that fell below 90\% were concentrated in
high-variance features (FIX, FVII, fibrinogen) where the real data exhibited large
inter-patient variability. PCA projections of the conditioned cohorts are provided
in Supplementary Figs.~S5--S6. These results demonstrated that SA's
multiplicity-weighted interpolation on the data manifold can meaningfully amplify
rare subgroups in a way that preserves their distinguishing clinical features---a
capability with direct relevance to hypothesis generation and power analysis in
understudied patient populations. MVN has no mechanism for this: fitting a separate
MVN to three PCOS patients across 216 features is mathematically impossible (rank~2
covariance), and post-hoc subsetting of unconditional MVN samples does not preserve
subgroup-specific structure.

\begin{figure}[ht]
\centering
\includegraphics[width=\textwidth]{validate_conditioned_features.pdf}
\caption{\textbf{Condition-specific feature preservation.} Grouped bar charts
comparing real (dark bars) and SA-generated synthetic (light bars) patient
means ($\pm$ SD error bars) for eight coagulation features, shown separately
for each clinical subgroup: Healthy ($n{=}14$, left, blue), PCOS ($n{=}3$,
center, orange), and PE ($n{=}5$, right, red). Values are pooled across all
three visits. Gray percentages indicate the mean relative error (MRE). Green
annotations show the bootstrap Mann--Whitney equivalence result: the fraction
of 1,000 bootstrap replicates (synthetic subsampled to $n_{\text{real}}$) in
which $p > 0.05$, i.e., real and synthetic could not be distinguished.
Features achieving $\geq$90\% are shown in green; $<$90\% in red. Across all
24 feature--condition pairs, 19 (79\%) were statistically indistinguishable in
$\geq$90\% of replicates (median: 96.2\%). The five features below this
threshold were concentrated in high-variance features (FIX, FVII, fibrinogen)
where large inter-patient variability reduced test power.}
\label{fig:conditioned-features}
\end{figure}

\subsection{Mechanistic Consistency}
\label{sec:results-l4}

As a final, orthogonal validation, we tested whether SA-generated patients produced
biologically plausible outputs when their coagulation factor levels were fed through
an independent mechanistic model of the coagulation cascade. We used the Hockin--Mann
BZ2012 model~\citep{Hockin2002,Luan2007}, a system of 58 ordinary differential
equations with 64 rate constants that simulates thrombin generation from
patient-specific coagulation factor inputs. We calibrated 5 of the 64 rate constants
on Visit~1 real patients at the population level (not per-patient), holding the
remaining 59 at published literature values. We then ran the calibrated model on all
23 real patients (3 visits $\times$ 2 conditions = 138 simulations, 0 failures) and
all 100 synthetic patients (3 visits $\times$ 2 conditions = 600 simulations, 0
failures), producing predicted thrombin generation features (lagtime, peak,
time-to-peak, maximum rate, and ETP) that we compared against the corresponding
measured values.

Under TF-only conditions, we found that synthetic patients produced the same
predicted-versus-measured cloud as real patients across all five TGA features
(Fig.~\ref{fig:mechanistic-scatter}). The median predicted-to-measured ratio was
0.92$\times$ for synthetic patients versus 0.91$\times$ for real patients, with
median absolute relative errors of 19.4\% and 21.1\%, respectively. We quantified
this overlap by computing the fraction of synthetic predicted/measured ratios that
fell within the 5th--95th percentile range of the real ratios and found cloud overlap
ranging from 86\% for lagtime to 94\% for maximum rate
(Fig.~\ref{fig:mechanistic-ratios}), which indicated that the SA-generated
factor-to-TGA mapping was indistinguishable from that of real patients at the
population level.

Under TF+TM conditions, the model systematically overcorrected the protein~C
anticoagulant pathway, producing peak and maximum rate predictions approximately
0.5$\times$ measured values for both real and synthetic patients (Supplementary
Fig.~S7). This shared systematic bias was itself informative: if synthetic patients
had violated the underlying coagulation biology, they would have deviated from real
patients in how the mechanistic model processed their factor levels. Instead, cloud
overlap under TF+TM remained high at 88--91\%.

We also verified that the calibration generalized across pregnancy timepoints. The
model was calibrated only on Visit~1 real patients, yet the predicted-to-measured
ratios at Visits~2 and~3 under TF-only conditions were 0.96$\times$ and 1.07$\times$,
respectively (Supplementary Fig.~S9), confirming that the calibrated rate constants
captured a stable population-level property of the coagulation system rather than
overfitting to the training visit. Spearman rank correlations between predicted and
measured values were moderate for ETP ($\rho \approx 0.6$--$0.8$ for real patients,
$\rho \approx 0.6$--$0.65$ for synthetic; Supplementary Fig.~S10) and weak for other
features, which was expected given that the Hockin--Mann model was not designed to
resolve inter-patient variability from a 5-parameter population-level calibration.
The key finding was not patient-level prediction but population-level plausibility:
an independent mechanistic model, encoding decades of coagulation biochemistry and
knowing nothing about the SA generation process, confirmed that synthetic patients
fell within the same biologically plausible envelope as real patients.

\begin{figure}[ht]
\centering
\includegraphics[width=\textwidth]{validate_mechanistic_tf_only.pdf}
\caption{\textbf{Mechanistic model predicted versus measured TGA features
(TF-only).} Each panel shows one TGA feature; the dashed line indicates perfect
prediction ($y{=}x$). Blue markers: real patients (circles = V1/train, squares
= V2, triangles = V3). Red/brown markers: SA-generated synthetic patients
(diamonds = V1, inverted triangles = V2, hexagons = V3). Synthetic patients
occupy the same predicted-versus-measured cloud as real patients across all
features and visits.}
\label{fig:mechanistic-scatter}
\end{figure}

\begin{figure}[ht]
\centering
\includegraphics[width=\textwidth]{validate_mechanistic_ratios_tf_only.pdf}
\caption{\textbf{Predicted/measured ratio distributions (TF-only).} Histograms
of the predicted-to-measured ratio for each TGA feature for real (blue) and
SA-generated synthetic (orange) patients. The dashed line indicates perfect
prediction (ratio = 1). The distributions overlap substantially, with cloud
overlap ranging from 86\% (lagtime) to 94\% (maximum rate).}
\label{fig:mechanistic-ratios}
\end{figure}
